\scanpage{5}
\uline{Théorème} 1.12 (Cartier) \deleted{Supposons que} \inserted{a) Si} dans $A$ les $\mathrm{Ker}$ de doubles flèches existent,$^{*}$ et \deleted{que} \inserted{si} $f$ y commute $(*)$, alors 1.11 est essentiellement surjectif,
% slanted
\margin{$*$ il suffit de regarder les doubles flèches $u,v$ telles que $\mathrm{Ker}(f(u),f(v))$ \uline{existe}}
\deleted{\unsure{et}} \inserted{b)} Pour que \uline{$h$} soit une équivalence [i.e. soit \deleted{de la forme} pleinement fidèle], il f. et s. que $f$ soit \uline{conservatif}, et que pour toute double flèche \inserted{$(u,v)$} de $A$ telle que $\mathrm{Ker}(f(u),f(v))$ existe, on ait que $\mathrm{Ker}(u,v)$ existe \uline{et} $f$ y commute.
% slanted
\margin{(condition \unsure{c|C})}
% slanted
\margin{améliorer a), en explicitant que $h$ est un \uline{passage au quotient}}

2) Partons d'une catégorie $B$, avec un foncteur
\[
 \varphi:B\mathrel{\unsure{-|\to}}B, \tag*{(2.1)}
\]
\deleted{satisfaisant les conditions (1.5),} et des liaisons \unsure{(1.2)}
\deleted{(2.1}
\[
 \varphi\mathrel{\overset{\alpha}{\unsure{-|\to}}}\mathrm{id}_B,\qquad\varphi\xrightarrow{\lambda}\varphi^2 \tag*{(2.2)}
\]
satisfaisant les conditions (1.5), (1.6). On peut alors (\(=\) sans ces restrictions (1.5) et (1.6) !) définir une catégorie
\[
 A=K(\varphi,\alpha,\lambda) \tag*{(2.3)}
\]
comme ci-dessus. Soit
\[
 A\xrightarrow{f}B \tag*{(2.4)}
\]
le foncteur \deleted{``de \unsure{réduction|oubli}'' \unsure{en}} canonique $(X,u)\longmapsto X$. On \deleted{\uline{Question} Quand $f$ admet-il un adjoint à droite $g$, de telle façon que \unsure{l'on ait} $\varphi\simeq fg$} va définir un \deleted{\unsure{objet}}foncteur
\[
 g:B\longrightarrow A \tag*{(2.5)}
\]
par
\[
 g(X)=\bigl(\varphi(X),\quad\varphi(X)\xrightarrow{\lambda(X)}\varphi^{2}(X)\bigr) \tag*{(2.6)}
\]
[Les conditions \inserted{\textcircled{\scriptsize 1.5}} \textcircled{\scriptsize (1.6)} assurent qu'on tombe bien dans $K(\varphi,\alpha,\lambda)$ !].
\deleted{Je dis qu'on a une relation d'adjonction,}
\[
 \deleted{\mathrm{Hom}_A\bigl((X,u),g(Y)\bigr)\simeq\mathrm{Hom}_B\bigl(\underset{\shortparallel\,X}{f(X,u)},\unsure{Y|X}\bigr)} \tag*{\deleted{(2.7)}}
\]
et on a évidemment
\[
 fg=\varphi \tag*{(2.7)}
\]
% archive sheet number: 4
